\chapter{NumPy による数値計算}
\label{chap:numpy}

NumPy\index{numpy@NumPy|textbf} は、Python で数値配列を扱うときの基本となるライブラリである。本章では、\code{ndarray} の形状と型、配列全体への演算、ブロードキャストを扱う。前半には配列の生成と変形、後半には集計、コピーとビュー、差分、ヒストグラム、線形代数の例を配置した。

\section{配列の構造とベクトル化}

\subsection{配列計算の基本概念}

Python のリストは柔軟だが、数値配列として大量の値を処理するには向いていないことがある。NumPy は、同じ型の数値を多次元配列 \code{ndarray}\index{ndarray@\code{ndarray}|textbf} として持ち、配列全体へまとめて演算をかけるためのライブラリである。最も基本的な NumPy 配列の定義と差分計算の例を示す。

\begin{lstlisting}[language=Python]
import numpy as np

times = np.array([10, 15, 23, 41], dtype=np.int64)
deltas = times[1:] - times[:-1]
\end{lstlisting}

この差分計算は Python の \code{for} ループでも記述できるが、NumPy では添字ループを明示せずに配列演算として記述できる。実行時間の比較方法は第~\ref{chap:performance}~章で扱う。図~\ref{fig:vectorization-benchmark} は比較図の模式例であり、この差分計算の実測結果ではない。

数式で書けば、これは隣り合う時刻の差
\[
\Delta t_i = t_{i+1} - t_i
\]
をまとめて計算していることに対応する。NumPy を使うと、このような「添字付きの式」を配列全体に対してそのまま書きやすい。

NumPy 配列は、Python のリストを高速化したものではない。同じ型の数値を共通のデータ領域に保持し、配列全体への演算を低水準の実装へ渡すデータ構造である。この構造に合わせて、要素ごとの手続きではなく、配列に適用する変換として処理を記述する。

\subsection{\code{ndarray} と属性}

NumPy の中心となるデータ構造は \code{ndarray} である。配列の生成後には、\code{shape}、\code{ndim}、\code{dtype} によって形状、次元数、要素型を確認する。

\begin{lstlisting}[language=Python]
import numpy as np

a = np.array([[1.0, 2.0, 3.0],
              [4.0, 5.0, 6.0]])

print(a.shape)     # (2, 3)
print(a.ndim)      # 2
print(a.dtype)     # float64
print(a.size)      # 6
print(a.itemsize)  # 8
\end{lstlisting}

\code{shape}\index{shape@\code{shape}|textbf} は各軸の長さ、\code{ndim} は次元数、\code{dtype}\index{dtype@\code{dtype}|textbf} は要素型、\code{size} は総要素数、\code{itemsize} は 1 要素あたりのバイト数である。値だけでなく \code{shape} と \code{dtype} も確認すると、形状不一致と型変換に起因する不具合を切り分けられる。

\begin{lstlisting}[numbers=none, xleftmargin=2\zw]
(2, 3)
2
float64
6
8
\end{lstlisting}

この出力は、\code{a} が 2 行 3 列、2 次元、6 要素の \code{float64} 配列であり、1 要素が 8 バイトであることを示す。形状不一致には \code{shape} と \code{ndim}、型変換には \code{dtype} と \code{itemsize} が診断材料となる。

\subsection{Python, C, NumPy でのデータ型の違い}
\label{sec:numpy-types}

Python の標準数値型は値ごとに型と管理情報を持つ。これに対して NumPy 配列では、配列全体に共通する固定長の型を \code{dtype} で指定する。代表的な数値型の対応を表~\ref{tab:type-comparison} に示す。

\begin{table}[hbtp]
\centering
\caption{Python, NumPy, C 言語での代表的な数値型の対応}
\label{tab:type-comparison}
\begin{tabular}{@{}p{0.14\linewidth}p{0.22\linewidth}p{0.22\linewidth}p{0.24\linewidth}@{}}
\toprule
分類 & C 言語の型 & Python の標準型 & NumPy の型 (\code{dtype}) \\
\midrule
整数型 & \code{int}, \code{long} & \code{int} (利用可能なメモリの範囲で可変長) & \code{int32}, \code{int64} \\
単精度浮動小数 & \code{float} (32ビット) & (標準では存在せず) & \code{float32} \\
倍精度浮動小数 & \code{double} (64ビット) & \code{float} & \code{float64} \\
論理型 & \code{bool} & \code{bool} & \code{bool\_} \\
\bottomrule
\end{tabular}
\end{table}

表のビット幅は一般的な処理系の例であり、C 言語の規格がすべての処理系に同じ幅を保証しているわけではない。C 言語では、\code{float} や \code{double} のように変数の型を宣言し、処理系が定める大きさの領域へ値を格納する。一般的な CPython では、標準の \code{int} は必要な桁数に応じて領域が増え、\code{float} は通常 C 言語の \code{double} と同じ倍精度形式を使う。また、各数値は型や参照管理の情報を含む Python オブジェクトである。多数の値を Python のリストへ格納すると、リスト中の参照と各オブジェクトの管理領域が必要になり、演算も要素ごとの Python 処理を経由する。

NumPy の \code{ndarray} は、要素ごとではなく配列全体で一つの \code{dtype} を共有し、データ領域と、形状・ストライドなどの管理情報を分けて保持する。配列は常に連続とは限らないが、同じ型の値を規則的にたどれるため、要素ごとの Python オブジェクト操作を避け、C や Fortran で実装されたループへまとめて処理を渡せる。この書き方をベクトル化と呼ぶ。

浮動小数点型は、必要な精度と値域、配列の大きさ、入出力元の型、CPU や GPU の特性を基準に選ぶ。NumPy の多くの配列生成関数や Python の \code{float} から作る配列では \code{float64} が既定になり、反復計算や差の小さい値を扱う計算では有力な選択肢になる。一方、\code{float32} は 1 要素が 4 バイトであり、8 バイトの \code{float64} と比べて配列本体のメモリ使用量が半分になる。大きな画像、センサデータ、GPU 計算、メモリ転送が支配的な処理では、この差が処理可能な配列サイズを左右する。元データが \code{float32} なら、必要性を検討せずに倍精度へ変換するとメモリ使用量だけが増える場合もある。

型を変えるときは、速度だけでなく結果が必要な精度を満たすかを確認する。例えば \code{np.finfo(np.float32).eps} と \code{np.finfo(np.float64).eps} を表示すると、1 付近で区別できる相対的な刻み幅を比較できる。どちらの型でも、桁落ちや丸め誤差そのものがなくなるわけではない。基準値が決まっていない段階では \code{float64} から始め、メモリや転送が問題になった時点で \code{float32} でも解析結果が変わらないかを検証する方法が堅実である。

\subsection{メモリ配置と低水準ループによる配列演算}

NumPy の配列演算と Python の要素別ループには、主に次の三つの実装上の差がある。

\begin{enumerate}
\item 数値を連続したメモリにまとめて持てること
\item Python レベルのループを減らせること
\item 内部で低レベルに最適化された処理を使えること
\end{enumerate}

配列演算は、Python で要素ごとに演算を呼び出す代わりに、配列全体を低水準の実装へ渡す。実行時間は、同じ入力と演算について Python ループと配列演算を測定して比較する。

配列演算では、差分、閾値判定、正規化、平均、分散を、明示的な添字ループを用いずに記述できる。処理対象となる配列と演算の対応が式に現れるため、要素ごとの更新処理と区別できる。

\section{配列の生成・演算・集計}

\subsection{配列の作成}

配列の生成時には、形状、要素型、初期値を定める。\code{np.array} は既存の値から配列を作り、\code{np.zeros} と \code{np.ones} は一定値で初期化し、\code{np.arange} と \code{np.linspace} は規則的な数列を作る。

\begin{lstlisting}[language=Python]
a = np.array([1, 2, 3], dtype=np.float64)
zeros = np.zeros((2, 3))
ones = np.ones((2, 3))
tmp = np.empty((2, 3))
grid = np.arange(0, 10, 2)
x = np.linspace(0.0, 1.0, 5)

print(a)
print(zeros.shape)
print(grid)
print(x)
\end{lstlisting}

\begin{lstlisting}[numbers=none, xleftmargin=2\zw]
[1. 2. 3.]
(2, 3)
[0 2 4 6 8]
[0.   0.25 0.5  0.75 1.  ]
\end{lstlisting}

\code{np.array}\index{numpy@NumPy!np.array@\code{np.array}|textbf} は既存の Python リストなどから配列を作る。\code{np.zeros}\index{numpy@NumPy!np.zeros@\code{np.zeros}} と \code{np.ones}\index{numpy@NumPy!np.ones@\code{np.ones}} は、すべての要素をそれぞれ 0 と 1 で初期化する。\code{np.arange}\index{numpy@NumPy!np.arange@\code{np.arange}} は指定した刻み幅の数列、\code{np.linspace}\index{numpy@NumPy!np.linspace@\code{np.linspace}} は指定区間を等分する数列を作る。後者はビン中心、時間軸、周波数軸の生成に利用できる。

\code{zeros.shape} が \code{(2, 3)} なので、\code{zeros} と \code{ones} は 2 行 3 列の配列である。\code{grid} は 0 から 8 まで 2 刻み、\code{x} は 0 から 1 を 5 点で等分した座標列である。\code{np.empty} は値を初期化せずにメモリを確保するため、生成直後の要素値は未定義である。読み出す前に全要素を代入する処理に用途を限定する。

\subsection{条件抽出・ソート・差分・ヒストグラム}

条件に合う要素の抽出、値のソート、隣接要素の差分、ヒストグラム\index{ひすとぐらむ@ヒストグラム|textbf}を、一つの時刻配列へ順に適用する例を示す。

\begin{lstlisting}[language=Python]
times = np.array([41, 10, 23, 15, 42])

selected = times[times > 20]
sorted_times = np.sort(times)
deltas = np.diff(sorted_times)
hist, edges = np.histogram(times, bins=[0, 20, 40, 60])

print(selected)
print(sorted_times)
print(deltas)
print(hist)
print(edges)
\end{lstlisting}

\begin{lstlisting}[numbers=none]
[41 23 42]
[10 15 23 41 42]
[ 5  8 18  1]
[2 1 2]
[ 0 20 40 60]
\end{lstlisting}

この例で行っていることは次のとおりである。

\begin{itemize}
\item フィルタ: \code{times > 20} は \code{[True, False, True, False, True]} に相当する条件配列を作り、その位置にある要素だけを \code{selected} として取り出す。
\item ソート: \code{np.sort(times)}\index{numpy@NumPy!np.sort@\code{np.sort}} は値を昇順に並べる。観測時刻やエネルギーのように順序が意味を持つ量では、後続の処理の前に並べ替えることが多い。
\item 差分: \code{np.diff(sorted\_times)}\index{numpy@NumPy!np.diff@\code{np.diff}} は隣り合う値の差を返す。返る配列の長さは元の配列より 1 つ短い。イベント時刻が昇順に並んでいれば、これは到着間隔として読める。
\item ヒストグラム: \code{np.histogram(times, bins=[0, 20, 40, 60])}\index{numpy@NumPy!np.histogram@\code{np.histogram}} は、区間ごとの個数を \code{hist}、区間境界を \code{edges} として返す。この例では \code{[0, 20)} に 2 個、\code{[20, 40)} に 1 個、\code{[40, 60]} に 2 個入っている。
\end{itemize}

条件式が返すブール配列は、そのまま添字に指定できる。複数条件には \code{(times > 20) \& (times < 40)} のように要素ごとの論理演算を用いる。\code{np.diff} の結果は要素の順序に依存するため、時刻差を求める前に入力の順序を確認し、必要ならソートする。

ヒストグラムは、ビン境界 \(b_0, b_1, \dots, b_n\) に対して
\[
h_k = \#\{ t_i \mid b_k \le t_i < b_{k+1} \}
\]
を数える操作と見なせる。ただし、NumPy では最後のビンだけ右端も含むので、最終区間は \([b_{n-1}, b_n]\) と読めばよい。

\subsection{要素ごとの演算とユニバーサル関数}

NumPy の演算子や多くの関数は、原則として要素ごとに働く。NumPy の公式ドキュメント~\cite{numpy_docs} では、\code{exp}、\code{sqrt}、\code{sin} のような要素ごと関数をユニバーサル関数\index{ゆにばーさるかんすう@ユニバーサル関数 (ufunc)|textbf}（ufunc）と呼んでいる。

\begin{lstlisting}[language=Python]
a = np.array([1.0, 2.0, 3.0])
b = np.array([10.0, 20.0, 30.0])

print(a + b)
print(a * 2)
print(a ** 2)
print(np.exp(a))
print(np.sqrt(a))
\end{lstlisting}

\begin{lstlisting}[numbers=none, xleftmargin=2\zw]
[11. 22. 33.]
[2. 4. 6.]
[1. 4. 9.]
[ 2.71828183  7.3890561  20.08553692]
[1.         1.41421356 1.73205081]
\end{lstlisting}

配列どうしの四則演算は、形がそろっていれば要素ごとに実行される。\code{np.exp(a)} や \code{np.sqrt(a)} も同様である。関数が 1 要素に対して何をするかを知っていれば、その関数が配列全体へどう作用するかも追いやすい。

たとえば、配列 \(a = (a_i)\)、\(b = (b_i)\) に対しては
\[
(a+b)_i = a_i + b_i,\qquad
(a^2)_i = a_i^2,\qquad
(\exp(a))_i = e^{a_i}
\]
のように、各要素へ同じ規則が適用されると考えればよい。

\subsection{集計と \code{axis}}

NumPy の集計関数は、配列全体または指定した軸に沿って値を集計する。\code{axis}\index{axis@\code{axis}|textbf} の指定と入出力の \code{shape} の関係を、次の例で確認する。

\begin{lstlisting}[language=Python]
matrix = np.array([[1.0, 2.0, 3.0],
                   [4.0, 5.0, 6.0]])

total = matrix.sum()
col_sum = matrix.sum(axis=0)
row_sum = matrix.sum(axis=1)
col_max = matrix.max(axis=0)

print(total)
print(col_sum)
print(row_sum)
print(col_max)
\end{lstlisting}

\begin{lstlisting}[numbers=none, xleftmargin=2\zw]
21.0
[5. 7. 9.]
[ 6. 15.]
[4. 5. 6.]
\end{lstlisting}

\code{total} は全要素の和である。\code{axis=0} は 0 軸に沿って集計するため各列の和を返し、\code{axis=1} は 1 軸に沿って集計するため各行の和を返す。入力と出力の \code{shape} を比較すると、集計によって除かれた軸を確認できる。

行列 \(A = (a_{ij})\) に対しては、
\[
(\mathrm{col\_sum})_j = \sum_i a_{ij},\qquad
(\mathrm{row\_sum})_i = \sum_j a_{ij}
\]
に対応する。したがって、\code{axis=0} の出力は列ごとの和、\code{axis=1} の出力は行ごとの和である。

\subsection{添字、スライス、条件抽出}

NumPy の添字には、Python リストと同じ整数・スライスに加え、多次元の軸指定とブール配列を使用できる。\index{すらいす@スライス|textbf}

\begin{lstlisting}[language=Python]
times = np.array([10, 15, 23, 41, 42])
print(times[0])
print(times[-1])
print(times[1:4])
print(times[times > 20])

grid = np.arange(12).reshape(3, 4)
print(grid[1, 2])
print(grid[:, 1])
print(grid[1, :])
\end{lstlisting}

\begin{lstlisting}[numbers=none, xleftmargin=2\zw]
10
42
[15 23 41]
[23 41 42]
6
[1 5 9]
[4 5 6 7]
\end{lstlisting}

\code{times[1:4]} は添字 1 から 3、すなわち 2 番目から 4 番目までを取り出す。\code{grid[:, 1]} は全行の 2 列目、\code{grid[1, :]} は 2 行目全体である。多次元配列に対する \code{for row in grid:} は、第 0 軸に沿って 1 行ずつ反復する。同じ演算を全行へ適用する場合は、明示的な反復と配列演算のどちらが入力と出力の関係を直接表すかを比較する。

\subsection{配列形状の変更}

配列の \code{shape} は、演算時の軸の対応を決める。1 次元配列、2 次元行列、2 次元の列ベクトルでは、同じ演算子を用いてもブロードキャスト後の結果が異なる。

\begin{lstlisting}[language=Python]
a = np.arange(12)
matrix = a.reshape(3, 4)
flat = matrix.ravel()
transposed = matrix.T
column = np.array([1.0, 2.0, 3.0])[:, None]

print(matrix.shape)
print(matrix)
print(flat)
print(transposed)
print(column.shape)
print(column)
\end{lstlisting}

\begin{lstlisting}[numbers=none, xleftmargin=2\zw]
(3, 4)
[[ 0  1  2  3]
 [ 4  5  6  7]
 [ 8  9 10 11]]
[ 0  1  2  3  4  5  6  7  8  9 10 11]
[[ 0  4  8]
 [ 1  5  9]
 [ 2  6 10]
 [ 3  7 11]]
(3, 1)
[[1.]
 [2.]
 [3.]]
\end{lstlisting}

\code{reshape} は配列の形状を変える。\code{ravel} は、可能であればビューのまま 1 次元配列へ変換する。\code{T} は転置を表す。\code{[:, None]} は新しい軸を 1 本追加し、1 次元配列を列ベクトルのように扱う書き方である。この軸の追加は、後述するブロードキャストで配列同士の対応関係を指定する役割を持つ。

\subsection{配列の連結と分割}

前処理済みデータを縦方向に連結する場合や、複数の特徴量を列方向に並べる場合には、連結する軸を指定する。逆に、一つの配列を同じ軸に沿って複数部分へ分割できる。

\begin{lstlisting}[language=Python]
x = np.array([1, 2, 3])
y = np.array([4, 5, 6])

joined = np.concatenate([x, y])
stacked = np.stack([x, y], axis=0)

matrix = np.arange(12).reshape(3, 4)
left, right = np.split(matrix, 2, axis=1)

print(joined)
print(stacked.shape)
print(stacked)
print(left)
print(right)
\end{lstlisting}

\begin{lstlisting}[numbers=none, xleftmargin=2\zw]
[1 2 3 4 5 6]
(2, 3)
[[1 2 3]
 [4 5 6]]
[[0 1]
 [4 5]
 [8 9]]
[[ 2  3]
 [ 6  7]
 [10 11]]
\end{lstlisting}

\code{concatenate} は既存の軸に沿って配列を連結し、\code{stack} は新しい軸を追加して配列を並べる。\code{split} は指定した軸に沿って配列を分割する。これらの操作では、実行後の \code{shape} が意図した値と一致するかを確認する。

\subsection{形とブロードキャスト}

NumPy では、配列の形状をそろえなくても、自動拡張の規則に従って演算できる場合がある。この規則をブロードキャスト\index{ぶろーどきゃすと@ブロードキャスト|textbf}という。

\begin{lstlisting}[language=Python]
x = np.array([1.0, 2.0, 3.0])
y = np.array([10.0, 20.0, 30.0])

z = x + y
matrix = x[:, None] + y[None, :]
print("z =")
print(z)
print("\nmatrix =")
print(matrix)
\end{lstlisting}

後半の \code{matrix} では、1 次元配列に軸を追加して、3 行 3 列の和の表を作っている。このように、形状が異なる配列を一定の規則で拡張して演算する仕組みをブロードキャストという。二つの 1 次元配列から要素の全組合せを計算できる一方、結果の要素数は両配列の要素数の積になる。この例の結果を次に示す。

\begin{lstlisting}[numbers=none, xleftmargin=2\zw]
z =
[11. 22. 33.]

matrix =
[[11. 21. 31.]
 [12. 22. 32.]
 [13. 23. 33.]]
\end{lstlisting}

この結果から、\code{z} は要素ごとの和、\code{matrix} は全組み合わせの表であることが分かる。\code{x} と \code{y} の要素数がそれぞれ \(N_x\) と \(N_y\) なら、\code{matrix} の要素数は \(N_xN_y\) になる。大きな配列では、生成前に形状、要素数、\code{dtype} から必要なメモリ量を見積もる。

数式で書けば、
\[
z_i = x_i + y_i,\qquad
M_{ij} = x_i + y_j
\]
である。後者の \(M_{ij}\) が、ブロードキャストによって作られる「全組み合わせの表」に対応している。

ブロードキャストでは形状を右端の軸から比較し、各軸の大きさが等しいか、一方が 1 なら演算できる。欠けている先頭の軸は大きさ 1 とみなす。たとえば \code{(2, 3)} と \code{(3,)} は適合するが、\code{(2, 3)} と \code{(2,)} は末尾の 3 と 2 が適合せず、\code{ValueError} になる。

\subsection{統計量とヒストグラム}

NumPy には、平均、分散、最小値、最大値、ヒストグラムなど、解析の基本となる集計関数がそろっている。代表的な統計量\index{numpy@NumPy!np.mean@\code{np.mean}}\index{numpy@NumPy!np.std@\code{np.std}}とヒストグラムを計算する例を下に示す。

\begin{lstlisting}[language=Python]
values = np.array([1.0, 2.0, 2.0, 3.0, 4.0])
mean = values.mean()
std = values.std(ddof=0)
hist, edges = np.histogram(values, bins=[0.0, 2.0, 4.0, 6.0])

print(f"mean = {mean:.2f}")
print(f"std = {std:.2f}")
print(hist)
print(edges)
\end{lstlisting}

\code{np.histogram} は、指定したビンごとの要素数とビン境界 \code{edges} を返し、図は作成しない。同じビン境界による複数データの比較や、度数に対するモデル計算を作図処理から独立させられる。これらの統計量を出力した結果を次に示す。

\begin{lstlisting}[numbers=none, xleftmargin=2\zw]
mean = 2.40
std = 1.02
[1 3 1]
[0. 2. 4. 6.]
\end{lstlisting}

この例では、\code{[0, 2)} に 1 個、\code{[2, 4)} に 3 個、\code{[4, 6]} に 1 個入っている。平均や標準偏差だけでは分布の形は分からず、ヒストグラムだけでは条件間比較がしにくい。数値要約と分布表示を組み合わせて読むことが、データ解析では基本になる。

\code{values.mean()} と \code{values.std(ddof=0)} は、それぞれ
\[
\bar{x} = \frac{1}{N}\sum_{i=1}^{N} x_i,\qquad
\sigma = \sqrt{\frac{1}{N}\sum_{i=1}^{N} (x_i-\bar{x})^2}
\]
を計算している。\code{ddof} は分母の自由度補正を指定し、一般に分母は \(N-\mathrm{ddof}\) となる~\cite{numpy_std}。\code{ddof=0} は手元の \(N\) 個の値のばらつきを分母 \(N\) で記述する指定であり、未知の母集団の標準偏差が求まったことを意味しない。独立に同じ分布から得た標本では、分母 \(N-1\) の標本分散が母分散の不偏推定量となる。ただし、その平方根が母標準偏差の不偏推定量になるわけではない。NumPy の \code{std()} の既定は \code{ddof=0}、pandas の \code{std()} の既定は \code{ddof=1} なので、比較時には指定をそろえる。
\index{ひょうじゅんへんさ@標準偏差|textbf}\index{ぶんさん@分散|textbf}\index{ddof@\code{ddof}|textbf}

\section{型・軸・ビューに関する注意点}

\subsection{演算子の意味の取り違え}

\code{*} は要素ごとの積であり、行列積ではない。行列積の計算には \code{@} または \code{np.dot} を用いる。

\begin{lstlisting}[language=Python]
A = np.array([[1.0, 2.0],
              [3.0, 4.0]])
B = np.array([[5.0, 6.0],
              [7.0, 8.0]])

print(A * B)
print(A @ B)
\end{lstlisting}

演算子を取り違えても、形状が適合すれば例外が発生せず、意味の異なる数値が得られる場合がある。線形代数や座標変換では、期待する数式と演算子を対応付け、既知の小さな入力で結果を検証する。

\subsection{\code{axis} と条件式}

2 次元配列を集計するとき、\code{axis=0} は第 0 軸を縮約して列ごとの値を返し、\code{axis=1} は第 1 軸を縮約して行ごとの値を返す。また、配列どうしの条件結合では \code{and} や \code{or} ではなく、要素ごとの論理演算子 \code{\&} や \code{|} を使う。

\begin{lstlisting}[language=Python]
matrix = np.array([[1, 2, 3],
                   [4, 5, 6]])
print(matrix.sum(axis=0))
print(matrix.sum(axis=1))

times = np.array([10, 15, 23, 41, 42])
mask = (times > 10) & (times < 40)
print(mask)
print(times[mask])
\end{lstlisting}

\begin{lstlisting}[numbers=none, xleftmargin=2\zw]
[5 7 9]
[ 6 15]
[False  True  True False False]
[15 23]
\end{lstlisting}

\code{times > 10 \& times < 40} のように括弧を省略すると、演算子の優先順位によって意図した比較にならない。各比較を括弧で囲み、その結果を \code{\&} または \code{|} で結合する。

\subsection{\code{dtype} と \code{np.empty}}

NumPy では、異なる型の値を同じ配列へ入れると、一定の変換規則に従って配列全体の \code{dtype} が決まる。その結果、数値計算を意図した配列が文字列配列になる場合がある。

\begin{lstlisting}[language=Python]
a = np.array([1, 2.5])
b = np.array([1, "2"])
tmp = np.empty(5)

print(a.dtype)
print(b.dtype.kind)
print(np.issubdtype(a.dtype, np.number))
print(np.issubdtype(b.dtype, np.number))
print(tmp.shape)
\end{lstlisting}

\begin{lstlisting}[numbers=none, xleftmargin=2\zw]
float64
U
True
False
(5,)
\end{lstlisting}

\code{a} は浮動小数点へそろえられる。これに対して \code{b} は \code{dtype.kind == "U"}、つまり Unicode 文字列の配列であり、数値演算には使えない。\code{tmp} は 0 で初期化されないため、読み出し前に全要素へ値を代入する。配列の生成方法に応じて、\code{dtype} と初期化の有無を確認する。

NumPy の整数型は固定長であり、表現範囲を超えるとオーバーフローし得る。特に符号なし整数の差は負の値を表現できないため、差分を求める前に、値域を確認して \code{astype(np.int64)} などで符号付き型へ変換する。Python の通常の \code{int} と同じ値域ではない。また、現在時刻付近の Unix 秒を \code{float64} にした値は 1 ns 刻みを保持できない。細かな時刻差が必要なら、範囲内の整数ナノ秒で差を取ってから秒へ換算する。

\subsection{コピーとビュー}

配列のコピーは独立したメモリを使用するため、大きな配列を繰り返しコピーするとメモリ使用量と転送時間が増える。一方、ビューは元配列とメモリを共有するため、ビューを通じた代入が元配列も変更する。各操作がコピーとビューのどちらを返すかを区別する必要がある。

\begin{lstlisting}[language=Python]
a = np.arange(6)

view = a[1:4]
view[:] = -1

copy = a[[0, 2, 4]]
copy[:] = 99

safe = a.copy()
print("a after view =", a)
print("copy         =", copy)
print("safe         =", safe)
\end{lstlisting}

この例では、スライスで作った \code{view} は元の \code{a} とデータを共有するので、\code{view[:] = -1} とすると \code{a} も変わる。これに対して、整数配列やブール配列による取り出しは新しい配列を返すので、\code{copy[:] = 99} としても元の \code{a} は変わらない。\code{reshape} や \code{ravel} も可能ならビューを返す。元配列から独立したデータ領域が必要な場合は、\code{copy()} を明示する。

\begin{lstlisting}[numbers=none, xleftmargin=2\zw]
a after view = [ 0 -1 -1 -1  4  5]
copy         = [99 99 99]
safe         = [ 0 -1 -1 -1  4  5]
\end{lstlisting}

この出力では、\code{view} を書き換えた時点で元配列 \code{a} が変化しているのに対し、\code{copy} は独立していることが見て取れる。

\subsection{浮動小数点数の比較}

浮動小数点数では、計算手順による丸め誤差のため、数学的に等しい値がビット単位では一致しない場合がある。許容する相対誤差と絶対誤差を定められる場合は、直接比較の代わりに \code{np.isclose} または \code{np.allclose} を用いる。

\begin{lstlisting}[language=Python]
a = 0.1 + 0.2
print(a == 0.3)
print(np.isclose(a, 0.3))
\end{lstlisting}

\begin{lstlisting}[numbers=none, xleftmargin=2\zw]
False
True
\end{lstlisting}

有限値に対する \code{np.isclose(a, b)} は、\(\lvert a-b\rvert\leq\mathrm{atol}+\mathrm{rtol}\lvert b\rvert\) を判定する。既定の許容差が測定精度に適するとは限らないため、単位と値の大きさに合わせて指定する。\code{np.allclose()} もブロードキャストするので、配列の同一性を検証する際は形状の一致を先に確認する。欠損値を一致とみなすかどうかも別に定める。

\section{差分・ヒストグラム・線形代数への応用}

\subsection{解析で NumPy を使うときの視点}

大規模データでは、必要以上に大きな型を選ぶとメモリ使用量が増える。たとえば、識別子や時刻差には \code{int64} が必要となる場合がある一方、範囲の狭いカテゴリ番号は \code{int16} や \code{int32} で表せる。要素数を \(N\)、1 要素のバイト数を \(b\) とすれば、配列本体には少なくとも \(Nb\) バイトが必要であるため、型の差は要素数に比例して総使用量へ反映される。

また、平均や分散だけでなく、差分、閾値判定、並べ替え、ヒストグラムまで NumPy の中で一貫して処理できると、コードの往復が減り、解析の流れを追いやすくなる。

\subsection{差分、ソート、ヒストグラム}

解析では、まず時刻順に並べ、次に差分を取り、その後で条件を掛けて分布を見る、という流れがよく現れる。NumPy ではこの流れをほぼそのまま式として書ける。

\begin{lstlisting}[language=Python]
times = np.array([10, 15, 23, 41, 42])
order = np.argsort(times)
sorted_times = times[order]
deltas = np.diff(sorted_times)
mask = (deltas >= 0) & (deltas <= 250)
hist, edges = np.histogram(deltas[mask], bins=50)
peak_bin = hist.argmax()
print("sorted_times =", sorted_times)
print("deltas       =", deltas)
print("peak_bin     =", peak_bin)
\end{lstlisting}

\code{np.argsort} は並べ替え順の添字を返す。\code{np.diff} は隣接差分である。\code{hist.argmax()} は最も高いビンの位置を返す。イベント時刻、スペクトル、波形特徴量のいずれでも、この種の処理は頻出である。

\begin{lstlisting}[numbers=none, xleftmargin=2\zw]
sorted_times = [10 15 23 41 42]
deltas       = [ 5  8 18  1]
peak_bin     = 0
\end{lstlisting}

\subsection{\code{numpy.linalg} による連立一次方程式の求解}

\code{numpy.linalg} は、連立一次方程式、行列分解、固有値などの線形代数演算を提供する。次の例では、2 元連立一次方程式 \(A\boldsymbol{x}=\boldsymbol{b}\) を \code{solve} で解く。

\begin{lstlisting}[language=Python]
A = np.array([[1.0, 1.0],
              [1.0, -1.0]])
b = np.array([3.0, 1.0])
x = np.linalg.solve(A, b)
\end{lstlisting}

行列演算は、校正係数の決定、座標変換、最小二乗法の前処理に現れる。制約付き最適化や疎行列など、\code{numpy.linalg} の範囲を越える機能は SciPy の公式資料~\cite{scipy_docs} を参照するとよい。

\subsubsection{Python リストで要素ごとに計算する場合}

Python の標準リストは異なる型のオブジェクトを同じリストに保持できる。一方、同じ数値演算を大量の要素へ適用する場合は、各要素について Python の演算を呼び出す必要がある。次の例では、100 万個の整数を 2 倍する処理を内包表記で記述している。

\begin{lstlisting}[language=Python]
# 100万要素のリストを作る
py_list = list(range(1000000))

# 内包表記の反復で、すべての要素を2倍にする
doubled_list = [x * 2 for x in py_list]
\end{lstlisting}

一般的な CPython のリストは、値そのものではなく Python オブジェクトへの参照を並べて保持する。内包表記では、各参照から値を取り出し、Python の整数演算を呼び出し、新しい整数オブジェクトへの参照を結果リストへ格納する。この反復処理が、固定長の数値配列を低レベルのループで処理する場合との差になる。

\subsubsection{NumPy 配列によるベクトル化}

NumPy 配列では、\code{int64} や \code{float64} のような共通の型に従って要素を格納する。\code{np\_array * 2} と書くと、Python で要素ごとのループを記述せず、NumPy 内部のループへ配列全体の処理を渡せる。リストを NumPy 配列へ変換してから計算する例を次に示す。

\begin{lstlisting}[language=Python]
import numpy as np

# 1. Python リストから共通 dtype の NumPy 配列を作る
np_array = np.array(py_list)

# 2. 配列全体の演算を NumPy の内部ループへ渡す
doubled_array = np_array * 2

# 3. 必要な出力形式がリストの場合にだけ変換する
back_to_list = doubled_array.tolist()
\end{lstlisting}

\subsubsection{pandas DataFrame との変換}

後述する pandas DataFrame は、列名と行ラベルを持つ表形式のデータ構造であり、列の内部では NumPy 配列や拡張配列が使われる。解析の途中で NumPy の関数を使いたい場合は、DataFrame や Series を NumPy 配列へ変換できる。pandas の DataFrame・Series と NumPy 配列を相互変換する基本形は次の通りである。

\begin{lstlisting}[language=Python]
import pandas as pd

# NumPy 配列から pandas の DataFrame（二次元表）に変換
df = pd.DataFrame(np_array, columns=["Value"])

# pandas の列から NumPy 配列を得る
values_array = df["Value"].to_numpy()
\end{lstlisting}

\code{to\_numpy()} がコピーを作るかどうかは、列の型、列数、必要な型変換に依存する。単一の数値列では元データとメモリを共有する場合もあるが、常に共有されるとは限らない。コピーの有無が性能や書き換えに影響する場面では、\code{np.shares\_memory()} で確認し、必要なら \code{to\_numpy(copy=True)} で独立した配列を要求する。

\subsection{公式資料の使い方}

NumPy の基礎を復習したいときは、NumPy quickstart~\cite{numpy_quickstart} を併読すると整理しやすい。特に、配列生成、添字とスライス、\code{shape} の変更、ブロードキャスト、コピーとビューは、解析コードを書くたびに戻って確認したくなる項目である。本章は解析で問題になりやすい点を補い、quickstart は関数名や細かな仕様を確認する資料と位置付けられる。

\section{章末課題：地震発生間隔の配列解析}

USGS の CSV をまだ用意していない場合は、次の配列を 1 問目と 2 問目の共通入力とせよ。この配列は秒単位の発生時刻であり、意図的に未整列にしてある。

\begin{lstlisting}[language=Python]
times = np.array([720.0, 0.0, 60.0, 4200.0, 180.0, 780.0])
\end{lstlisting}

\begin{enumerate}
\item \textbf{発生時刻からの発生間隔配列の作成}。
第~\ref{chap:io}~章で扱った CSV から \code{time} 列を読み、秒単位の数値配列へ変換せよ。その後で NumPy 配列へまとめ、\code{np.sort} と \code{np.diff} を使って隣り合う地震の発生間隔 \(\Delta t\) を求めよ。
小規模入力では、整列後の時刻が \code{[0., 60., 180., 720., 780., 4200.]}、\(\Delta t\) が \code{[60., 120., 540., 60., 3420.]} になる。\code{np.sort()} が返す配列と \code{np.diff()} が返す配列について、\code{shape}、\code{dtype}、要素数が 1 減る理由を説明せよ。

\item \textbf{NumPy による発生間隔の基本統計とヒストグラム}。
\(\Delta t\) 配列に対して、平均値、0 秒超だけに絞った要素数、\([0, 60)\,\mathrm{s}\)、\([60, 600)\,\mathrm{s}\)、\([600, 3600)\,\mathrm{s}\) の 3 区間の個数を求めよ。さらに、\code{edges = np.arange(0.0, 3600.0 + 600.0, 600.0)} として 600 秒幅の境界を作り、\code{np.histogram} によって各ビンの個数と最頻ビンの範囲を求めよ。
小規模入力の平均値は 840 秒であり、0 秒超の要素数は 5 である。3 区間の個数は順に 0、4、1 になる。600 秒幅のヒストグラムでは、\([0,600)\,\mathrm{s}\) が 4 件で最頻ビンとなり、最後の \([3000,3600]\,\mathrm{s}\) が 1 件となる。区間の右端を含めるかどうかが条件式と一致していることを確認し、\code{np.histogram()} が返す \code{counts} と \code{bin\_edges} の長さの関係も説明せよ。

\item \textbf{発展課題}。
発生間隔は短時間側に密集し、長時間側には少数の値が残る場合がある。そのため、等間隔ビンだけでは短時間側の差を分解できないことがある。たとえば \code{np.logspace(np.log10(1), np.log10(3600), 30)} のようにビン境界を作り、\(\Delta t\) のヒストグラムを対数ビンで作り直せ。固定幅ビンと対数ビンについて、ビン幅、空のビンの数、最短値から最長値までを収めるビン数を比較せよ。発展課題では、日本近海の範囲を緯度・経度で定義してデータを別に取得し、同じ集計条件による全球データとの分布差を比較せよ。
\end{enumerate}

NumPy 配列は、配列全体で一つの \code{dtype} を共有する。一方、観測時刻、検出器名、信号値のように型の異なる列を一つの記録として保持するには、列ごとに型を持つ表形式の構造を用いる。次章では、pandas の列名、欠損値、グループ集計、結合を扱う。
